% Part: lambda-calculus
% Chapter: church-rosser
% Section: parallel-beta-reduction

\documentclass[../../../include/open-logic-section]{subfiles}

\begin{document}

\olfileid[fr]{lam}{cr}{pb}
\olsection{Réduction~\texorpdfstring{$\beta$}{β} parallèle}

Nous introduisons la \emph{réduction~$\beta$
parallèle} et montrons qu'elle vérifie la
propriété de Church--Rosser. Les règles portent
sur les classes d'$\alpha$-équivalence. On peut
renommer les lieurs dans les représentants.
La substitution utilisée dans la règle de
contraction est celle du \olref[lam][syn][alp]{cor:sub}.

\begin{defn}[Réduction~$\beta$ parallèle, $\bredpar$]
\ollabel{defn:bredpar}
La réduction parallèle ($\bredpar$) des
termes est définie inductivement par les
règles suivantes :
\begin{enumerate}
\item \ollabel{defn:bredpar1} $x \bredpar x$.
\item \ollabel{defn:bredpar2}
  Si $N \bredpar N'$, alors
  $\lambd[x][N] \bredpar \lambd[x][N']$.
\item \ollabel{defn:bredpar3}
  Si $P \bredpar P'$ et $Q \bredpar Q'$,
  alors $PQ \bredpar P'Q'$.
\item \ollabel{defn:bredpar4}
  Si $N \bredpar N'$ et $Q \bredpar Q'$,
  alors $(\lambd[x][N])Q \bredpar
  \Subst{N'}{Q'}{x}$.
\end{enumerate}
\end{defn}

La règle pour l'abstraction permet de conserver le corps
de celle-ci. Par exemple, la première règle donne
$x\bredpar x$ et la deuxième donne alors
$\lambd[x][x]\bredpar\lambd[x][x]$.

La réduction~$\beta$ parallèle permet de
contracter en une étape un nombre quelconque
de rédex présents dans le terme initial.
Elle ne contracte pas dans cette même étape
un rédex qui vient seulement d'être créé.
Ainsi, $(\lambd[f][fx])(\lambd[y][y])$
se réduit parallèlement à lui-même ou à
$(\lambd[y][y])x$, mais pas directement
à~$x$. Le rédex dont la contraction donne
$x$ n'apparaît qu'après la première étape;
une seconde étape parallèle suffit.
La réduction~$\beta$ ordinaire, avec
plusieurs étapes, conduit bien à~$x$.

\begin{thm}\ollabel{thm:refl}
$M \bredpar M$.
\end{thm}

\begin{proof}
Exercice.
\end{proof}

\begin{prob}
Démontrer le \olref[lam][cr][pb]{thm:refl}.
\end{prob}

\begin{defn}[Développement~$\beta$ complet]
\ollabel{defn:bcd}
Le \emph{développement~$\beta$ complet}
$\bcd{M}$ de $M$ est défini inductivement
par les clauses suivantes :
\begin{align}
  \bcd{x} &= x \ollabel{defn:bcd1} \\
  \bcd{(\lambd[x][N])}
    &= \lambd[x][\bcd{N}] \ollabel{defn:bcd2}\\
  \bcd{(PQ)} &= \bcd{P}\bcd{Q}
    && \text{si $P$ n'est pas une abstraction $\lambd$}
    \ollabel{defn:bcd3} \\
  \bcd{((\lambd[x][N])Q)}
    &= \Subst{\bcd{N}}{\bcd{Q}}{x}
    \ollabel{defn:bcd4}
\end{align}
\end{defn}

Comme son nom l'indique, le développement
$\beta$ complet contracte en parallèle
tous les rédex du terme initial. La
réduction parallèle permet au contraire
d'en laisser certains intacts. Ainsi, le
développement complet de
$(\lambd[f][fx])(\lambd[y][y])$ est
$(\lambd[y][y])x$, et non le terme initial.

Chaque appel récursif de la définition de $\bcd{M}$
porte sur un sous-terme strict de $M$, dont le nombre
de nœuds est inférieur. La dernière clause applique
la substitution aux développements déjà obtenus.
Elle se lit à $\alpha$-équivalence près, avec un
représentant qui rend la substitution définie.

\begin{lem}\ollabel{lem:comp}
Si $M \bredpar M'$ et $R \bredpar R'$,
alors $\Subst{M}{R}{y} \bredpar
\Subst{M'}{R'}{y}$.
\end{lem}

\begin{proof}
Pour composer les substitutions, on utilise
l'identité suivante, à $\alpha$-équivalence près.
Si $x\neq y$ et $x\notin\FV{V}$, alors
\[
  \Subst{\Subst{U}{W}{x}}{V}{y}
  \aeq
  \Subst{\Subst{U}{V}{y}}{\Subst{W}{V}{y}}{x}.
\]
Pour le voir, choisissons dans $U$ et $W$ des noms de
lieurs distincts de $x$, de $y$ et de toutes les variables
libres de $U$, $W$ et $V$. Dans l'arbre étiqueté de $U$
utilisé dans la preuve du \olref[lam][syn][alp]{thm:sub},
le membre gauche insère $W$ aux feuilles libres de nom $x$,
puis $V$ aux feuilles libres de nom $y$, y compris celles
des copies insérées de $W$. Le membre droit remplace
d'abord les feuilles libres de nom $y$ par $V$ et insère
ensuite, aux feuilles libres de nom $x$, des copies de
$\Subst{W}{V}{y}$. Aucune copie de $V$ n'ajoute de feuille
libre de nom $x$. Les deux opérations produisent donc le
même arbre étiqueté.

Raisonnons par induction sur la dérivation
de $M \bredpar M'$.
\begin{enumerate}
\item Si la dernière étape applique la règle
  \olref{defn:bredpar1}, on a $M=M'=z$ pour une
  variable $z$. Si $z=y$, la conclusion est
  $R\bredpar R'$. Si $z\neq y$, c'est
  $z\bredpar z$.
\item Si la dernière étape applique la règle
  \olref{defn:bredpar2}, alors
  $M=\lambd[x][N]$ et
  $M'=\lambd[x][N']$, avec
  $N\bredpar N'$. En prenant un représentant
  dont $x$ est frais pour $y$, $R$ et $R'$,
  l'objectif devient
  $\lambd[x][\Subst{N}{R}{y}]
  \bredpar
  \lambd[x][\Subst{N'}{R'}{y}]$.
  Il découle de l'hypothèse d'induction
  et de la règle \olref{defn:bredpar2}.
\item Si la dernière étape applique la règle
  \olref{defn:bredpar3}, les hypothèses d'induction
  s'appliquent aux deux termes de l'application.
  La même règle réunit les deux réductions obtenues.
\item Si la dernière étape applique la règle
  \olref{defn:bredpar4}, alors
  $M=(\lambd[x][N])Q$ et
  $M'=\Subst{N'}{Q'}{x}$.
  Prenons encore $x$ frais pour $y$, $R$
  et $R'$. Les hypothèses d'induction
  donnent
  $\Subst{N}{R}{y}\bredpar
  \Subst{N'}{R'}{y}$ et
  $\Subst{Q}{R}{y}\bredpar
  \Subst{Q'}{R'}{y}$.
  La règle \olref{defn:bredpar4} donne
  alors
  \[
    (\lambd[x][\Subst{N}{R}{y}])
    \Subst{Q}{R}{y}
    \bredpar
    \Subst{\Subst{N'}{R'}{y}}
          {\Subst{Q'}{R'}{y}}{x}.
  \]
  Par composition, ce dernier terme est
  $\alpha$-équivalent à
  $\Subst{\Subst{N'}{Q'}{x}}{R'}{y}$,
  comme le montrent les conditions $x\neq y$
  et $x\notin\FV{R'}$ vérifiées par le choix de~$x$.
\end{enumerate}
\end{proof}

\begin{prob}
Dans le \olref[lam][cr][pb]{lem:comp}, écrire
explicitement les deux applications de l'hypothèse
d'induction lorsque la dernière règle porte sur une
application. Vérifier le rôle de $x\notin\FV{V}$ dans
l'identité de composition en prenant $U=y$.
\end{prob}

\begin{lem}\ollabel{lem:cont}
Si $M \bredpar M'$, alors
$M' \bredpar \bcd{M}$.
\end{lem}

\begin{proof}
Raisonnons par induction sur la dérivation
de $M \bredpar M'$.
\begin{enumerate}
\item Si la dernière étape applique la règle
  \olref{defn:bredpar1}, c'est un exercice.
\item Si la dernière étape applique la règle
  \olref{defn:bredpar2}, alors
  $M=\lambd[x][N]$ et
  $M'=\lambd[x][N']$, avec
  $N\bredpar N'$. Par l'hypothèse
  d'induction et la règle \olref{defn:bredpar2},
  $\lambd[x][N']\bredpar
  \lambd[x][\bcd{N}]$; la cible est
  $\bcd{(\lambd[x][N])}$ par
  l'\olref{defn:bcd2}.
\item Si la dernière étape applique la règle
  \olref{defn:bredpar3}, alors
  $M=PQ$ et $M'=P'Q'$, où
  $P\bredpar P'$ et $Q\bredpar Q'$.
  Les hypothèses d'induction donnent
  $P'\bredpar\bcd{P}$ et
  $Q'\bredpar\bcd{Q}$.
  \begin{enumerate}
  \item Si $P=\lambd[x][N]$, alors
    $P'=\lambd[x][N']$ pour un $N'$
    tel que $N\bredpar N'$. Puisque
    $N'\bredpar\bcd{N}$ et
    $Q'\bredpar\bcd{Q}$, la règle
    \olref{defn:bredpar4} donne
    $(\lambd[x][N'])Q'\bredpar
    \Subst{\bcd{N}}{\bcd{Q}}{x}$.
  \item Si $P$ n'est pas une abstraction
    $\lambd$, la règle
    \olref{defn:bredpar3} donne
    $P'Q'\bredpar\bcd{P}\bcd{Q}$,
    qui est $\bcd{PQ}$ par
    l'\olref{defn:bcd3}.
  \end{enumerate}
\item Si la dernière étape applique la règle
  \olref{defn:bredpar4}, alors
  $M=(\lambd[x][N])Q$ et
  $M'=\Subst{N'}{Q'}{x}$, avec
  $N\bredpar N'$ et $Q\bredpar Q'$.
  Par l'hypothèse d'induction,
  $N'\bredpar\bcd{N}$ et
  $Q'\bredpar\bcd{Q}$. Le
  \olref{lem:comp} donne
  $\Subst{N'}{Q'}{x}\bredpar
  \Subst{\bcd{N}}{\bcd{Q}}{x}$;
  le terme de droite est précisément
  $\bcd{((\lambd[x][N])Q)}$.
\end{enumerate}
\end{proof}

\begin{prob}
Compléter la preuve du \olref[lam][cr][pb]{lem:cont}.
\end{prob}

\begin{thm}\ollabel{thm:cr}
$\bredpar$ vérifie la propriété de
Church--Rosser.
\end{thm}

\begin{proof}
Si $M\bredpar P$ et $M\bredpar Q$,
le \olref{lem:cont} donne
$P\bredpar\bcd{M}$ et
$Q\bredpar\bcd{M}$. Ce terme est le
sommet commun requis.
\end{proof}

\end{document}
